\section*{Hoofdstuk \ref{ch:g9:gravitation} --- Gravitatie}

\begin{solution}{exo:g9:gravitation:1}
$F = G \, m_A m_B / d^2$, met de \omterm{def:g3:measuring-mass:mass}{massa}'s in
\omterm{def:g3:measuring-mass:units}{kilogram}, $d$ in \omterm{def:g2:measuring-length:units}{meter} tussen de middelpunten,
$F$ in newton en $G = \qty{6.67e-11}{N.m^2/kg^2}$. In woorden: de trek
groeit met elke \omterm{def:g3:measuring-mass:mass}{massa} en sterft met het kwadraat van de afstand
--- en ze werkt even sterk op beide partners, in tegengestelde
richtingen.
\end{solution}

\begin{solution}{exo:g9:gravitation:2}
Met precies \qty{1}{N} --- de trekken van de wet komen altijd in even
grote, tegengestelde paren. De appel versnelt zichtbaar omdat zijn
\omterm{def:g3:measuring-mass:mass}{massa} piepklein is; dezelfde \qty{1}{N} op \qty{5.97e24}{kg}
\omterm{def:g4:solar-system:planet}{planeet} beweegt die onmeetbaar weinig.
\end{solution}

\begin{solution}{exo:g9:gravitation:3}
Op $2d$: een kwart van de \omterm{def:g2:pushes-and-pulls:force}{kracht} op $d$ --- het kwadraat van
twee. Op $10d$: een honderdste.
\end{solution}

\begin{solution}{exo:g9:gravitation:4}
$F = \num{6.67e-11} \times 80 \times 80 \div 0.25 = \qty{1.7e-6}{N}$
--- minder dan twee miljoensten van een newton: wat danspartners ook
naar elkaar toe trekt, \omterm{def:g9:gravitation:gravitation}{gravitatie} is het niet.
\end{solution}

\begin{solution}{exo:g9:gravitation:5}
De \omterm{def:g7:motion-average-speed:formula}{snelheid} waarmee de kogel valt, en de \omterm{def:g7:motion-average-speed:formula}{snelheid} waarmee het oppervlak
van de ronde Aarde eronder wegkrult. Een \omterm{def:g6:describing-motion:trajectory}{baan}: vallen en
missen.
\end{solution}

\begin{solution}{exo:g9:gravitation:6}
Op de hoogte van het station behoudt de zwaartekracht bijna haar volle
sterkte aan het oppervlak --- zij is juist wat het station op zijn
cirkel houdt. Station en bemanning zijn allebei aan het \emph{vallen},
samen, om de Aarde heen: metgezellen die samen vallen drukken nergens
op, en zweven is het gevoel van gedeelde vrije val.
\end{solution}

\begin{solution}{exo:g9:gravitation:7}
Omdat de hemel in monsterlijke \omterm{def:g3:measuring-mass:mass}{massa}'s handelt ---
planeten en sterren vermenigvuldigen de zwakke $G$
tot machtige \omterm{def:g2:pushes-and-pulls:force}{krachten} --- en omdat
\omterm{def:g9:gravitation:gravitation}{gravitatie} nooit wordt opgeheven: ze
trekt alleen maar aan, dus verdiept elke
toegevoegde \omterm{def:g3:measuring-mass:mass}{massa} de trek.
\end{solution}

\begin{solution}{exo:g9:gravitation:8}
De trek van de \omterm{def:g2:moon-in-the-sky:moon}{Maan} stapelt de zeeën zowel naar haar toe (het
dichtstbijzijnde water, het hardst getrokken) als van haar af (het
verste water, het minst getrokken); de draaiende Aarde voert elke kust
dagelijks onder beide bulten door --- twee getijden. Het grootst bij
nieuwe en \omterm{ex:g2:moon-in-the-sky:shapes}{volle maan}, wanneer de hand van de Zon op één lijn met
die van de \omterm{def:g2:moon-in-the-sky:moon}{Maan} staat: springtij.
\end{solution}

\begin{solution}{exo:g9:gravitation:9}
$F = \num{6.67e-11} \times (\num{7.3e22} \times \num{5.97e24}) \div
(\num{3.8e8})^2 \approx \qty{2.0e20}{N}$ --- een $2$ gevolgd door
twintig nullen aan newton die de \omterm{def:g2:moon-in-the-sky:moon}{Maan} aan haar pacht houden.
\end{solution}

\begin{solution}{exo:g9:gravitation:10}
Twee knoppen draaien tegen elkaar in: de \omterm{def:g3:measuring-mass:mass}{massa} van de
\omterm{def:g2:moon-in-the-sky:moon}{Maan}, ongeveer tachtig keer kleiner, deelt de trek door tachtig;
maar haar straal, bijna vier keer kleiner, \emph{versterkt} de trek aan
het oppervlak met vier in het kwadraat --- ongeveer veertien keer.
Tachtig omlaag, veertien omhoog: er blijft ongeveer een zesde over.
\end{solution}

\begin{solution}{exo:g9:gravitation:11}
Het omgekeerde kwadraat verdunt de greep van de Zon met de afstand, en
een zwakkere greep kan alleen een tragere val-eromheen vasthouden: de
buitenste planeten lopen los aan de lijn en kuieren, de binnenste
worden stevig vastgehouden en racen --- de sprint van Mercurius en de
wandeling van 164 jaar van Neptunus zijn dezelfde wet, op twee
afstanden gelezen.
\end{solution}

\begin{solution}{exo:g9:gravitation:12}
Zijn omlooptijd moet precies één dag zijn, en hij moet naar het oosten
cirkelen --- de draairichting van de Aarde zelf. Dan draaien satelliet
en antenne in volmaakte pas: de ``hangende'' schotel kijkt naar een
kanonskogel die precies zo snel om de \omterm{def:g4:solar-system:planet}{planeet} heen valt als de
\omterm{def:g4:solar-system:planet}{planeet} eronder ronddraait --- bewegingloos alleen in de ogen van
de meedraaiende waarnemer.
\end{solution}

\begin{solution}{pb:g9:gravitation:1}
\textbf{1.} $F = \num{6.67e-11} \times (0.1 \times \num{5.97e24}) \div
(\num{6.4e6})^2 \approx \qty{0.97}{N}$ --- noem het één newton: de ruk
uit de boomgaard, berekend uit de constante van het heelal.
\textbf{2.} Eveneens \qty{0.97}{N} --- wederkerigheid: de trekken van de
wet komen alleen in even grote, tegengestelde paren.
\textbf{3.} Het verdubbelen van de afstand tussen de middelpunten
deelt de trek door vier: ongeveer \qty{0.24}{N}.
\textbf{4.} Elke kruimel van de \omterm{def:g4:solar-system:planet}{planeet} trekt aan de appel ---
bergen dichtbij, de kern diep, de tegenvoeters ver --- en de som van al
die trekken pakt precies zo uit alsof de hele \omterm{def:g3:measuring-mass:mass}{massa} in het
middelpunt zat: een stelling, geen aanname, en de auteur ervan stelde de
publicatie jaren uit tot hij haar kon bewijzen.
\textbf{5.} Ongeveer $2000 \times 9.8 \approx \qty{2.0e4}{N}$ ---
twintigduizend newton (de wet met $d$ gelijk aan één aardstraal geeft
hetzelfde).
\textbf{6.} Op twee stralen laat het omgekeerde kwadraat een kwart
over: ongeveer \qty{4900}{N} tegenover zijn \qty{19600}{N} aan
\omterm{def:g4:weight-and-mass:weight}{gewicht} aan het oppervlak.
\textbf{7.} Ja --- een kwart van de volle zwaartekracht is geen
gewichtloze leegte. De satelliet geeft die trek uit aan het krommen van
zijn \omterm{def:g6:describing-motion:trajectory}{baan}: hij valt om de Aarde heen en mist eeuwig, precies zoals het
bergkanon het leerde.
\textbf{8.} ``Snel genoeg afgevuurd past de val van een kanonskogel bij
de kromming van de Aarde, en cirkelt hij voor altijd --- vallend zonder
te landen. Elke satelliet is die kanonskogel, zijwaarts snel genoeg
gelanceerd om de grond te blijven missen.''
\textbf{9.} $\num{3.8e8} \div \num{6.4e6} \approx 59.4$ aardstralen.
\textbf{10.} Ongeveer $1 \div 59.4^2 \approx 1/3500$ van de sterkte aan
het oppervlak.
\textbf{11.} Zacht omdat de zwaartekracht op zestig stralen
drieënhalfduizendvoudig verdund is --- de val van
vijf \omterm{def:g2:measuring-length:units}{meter} in de eerste \omterm{def:g2:measuring-time:second}{seconde} krimpt tot ongeveer een
millimeter. Eindeloos omdat die millimeter val precies past bij de
millimeter waarmee het oppervlak van de Aarde onder de zijwaartse glijvlucht
van de \omterm{def:g2:moon-in-the-sky:moon}{Maan} wegkrult: de val loopt nooit in.
\textbf{12.} Bijvoorbeeld: ``Appel, satelliet en \omterm{def:g2:moon-in-the-sky:moon}{Maan}
gehoorzamen één zin: elke \omterm{def:g3:measuring-mass:mass}{massa} valt naar elke andere toe,
volgens $G m_A m_B / d^2$ --- de appel landt, de satelliet en de
\omterm{def:g2:moon-in-the-sky:moon}{Maan} blijven missen, en het verschil is alleen zijwaartse
\omterm{def:g7:motion-average-speed:formula}{snelheid}.''
\end{solution}
